\documentclass{report}
\usepackage[utf8]{inputenc}
\usepackage[T1]{fontenc}
\usepackage[french]{babel}
\usepackage{amsmath, amssymb, amsthm}
\usepackage{graphicx}
\usepackage{xcolor}
\usepackage{listings}
\usepackage{tikz}
\usepackage{hyperref}
\usepackage{geometry}
\geometry{a4paper, margin=2cm}

\title{Suites de fonctions, étude de la convergence simple et de la convergence uniforme}
\author{Charles EDOU NZE}
\date{2026-07-18}

\lstset{
    language=Python,
    basicstyle=\ttfamily\small,
    keywordstyle=\color{blue},
    stringstyle=\color{red},
    commentstyle=\color{green!60!black},
    showstringspaces=false,
    frame=single,
    breaklines=true,
    literate={é}{{\'e}}1 {è}{{\`e}}1 {à}{{\`a}}1 {ç}{{\c c}}1 {ù}{{\`u}}1 {ê}{{\^e}}1 {î}{{\^i}}1 {ô}{{\^o}}1 {û}{{\^u}}1 {â}{{\^a}}1 {É}{{\'E}}1 {È}{{\`E}}1 {À}{{\`A}}1 {Ç}{{\c C}}1 {Ù}{{\`U}}1 {Ê}{{\^E}}1 {Î}{{\^I}}1 {Ô}{{\^O}}1 {Û}{{\^U}}1 {Â}{{\^A}}1
}

\begin{document}
\maketitle
\tableofcontents
\newpage


\chapter*{Jalon 21 : Suites de fonctions, étude de la convergence simple et de la convergence uniforme}

\section*{1. Présentation du concept clé : La danse des approximations}

\begin{figure}[h]
\centering
\begin{tikzpicture}[scale=4]
  % Axes
  \draw[->] (-0.1,0) -- (1.2,0) node[right] {$x$};
  \draw[->] (0,-0.1) -- (0,1.2) node[above] {$f_n(x)$};

  % Grid
  \draw[dotted] (1,0) -- (1,1);
  \draw[dotted] (0,1) -- (1,1);

  % Functions x^n
  \draw[thick, blue, domain=0:1, samples=100] plot (\x, {\x});
  \node[blue, right] at (0.6, 0.7) {$n=1$};

  \draw[thick, red, domain=0:1, samples=100] plot (\x, {\x^2});
  \node[red, right] at (0.7, 0.4) {$n=2$};

  \draw[thick, green!70!black, domain=0:1, samples=100] plot (\x, {\x^5});
  \node[green!70!black, right] at (0.8, 0.2) {$n=5$};

  \draw[thick, orange, domain=0:1, samples=100] plot (\x, {pow(\x, 20)});
  \node[orange, right] at (0.9, 0.05) {$n=20$};

  % Limit function
  \draw[ultra thick, black, dashed] (0,0) -- (0.98,0);
  \draw[ultra thick, black, dashed] (0.98,0) -- (1,1);
  \fill[black] (1,1) circle (0.5pt);

  \node at (0.5, -0.15) {Convergence simple, mais non uniforme};
\end{tikzpicture}
\caption{Suite de fonctions $f_n(x) = x^n$ sur $[0,1]$ convergeant vers la fonction limite discontinue.}
\end{figure}


Imaginez que vous êtes un observateur attentif d'un phénomène complexe, par exemple l'évolution de la température le long d'une barre métallique chauffée. Chaque jour, vous effectuez des mesures, et ces mesures, prises à différents points de la barre, décrivent une certaine courbe de température. Appelons ces courbes $f_1(x)$, $f_2(x)$, $f_3(x)$, et ainsi de suite, où $x$ est la position le long de la barre et l'indice représente le jour. Vous obtenez ainsi une "suite" de fonctions.

L'objectif est de comprendre si cette suite de fonctions "tend" vers une situation finale stable, une sorte de profil de température d'équilibre. Mais qu'est-ce que cela signifie, pour une suite de fonctions, de tendre vers une autre fonction ?

Considérons une première approche : chaque jour, pour chaque point $x$ de la barre, la température $f_n(x)$ se rapproche de plus en plus de la température d'équilibre $f(x)$. C'est une convergence "point par point". Si je me fixe sur un point précis de la barre, disons son extrémité droite, et que je regarde la suite des températures à cet endroit ($f_1(x_0), f_2(x_0), \dots$), cette suite de nombres converge vers $f(x_0)$. Si cela est vrai pour \textit{tous} les points de la barre, on parle de \textbf{convergence simple}. C'est une forme de rapprochement, certes, mais elle peut être trompeuse.

Imaginons une rangée de musiciens qui accordent leurs instruments. La convergence simple, c'est comme si chaque musicien réussissait à accorder parfaitement son instrument. L'un après l'autre, ils trouvent la note juste. Cependant, il n'y a aucune garantie qu'ils soient tous accordés \textit{en même temps}. Il se pourrait que lorsque le premier a fini, le dernier est encore loin du compte, et inversement. Le temps nécessaire pour que chaque musicien atteigne l'accord parfait pourrait dépendre du musicien lui-même.

Maintenant, visualisez une seconde approche : non seulement chaque point de la barre voit sa température se stabiliser, mais l'écart maximum entre le profil de température du jour $n$ et le profil d'équilibre diminue globalement. Autrement dit, le "pire" désaccord sur l'ensemble de la barre devient de plus en plus petit à mesure que les jours passent. Ici, le rythme d'accordement n'est pas point par point, mais global. On exige qu'à partir d'un certain jour $N$, l'ensemble de la barre présente un écart de température partout inférieur à une petite marge $\epsilon$. C'est ce que l'on nomme la \textbf{convergence uniforme}.

Pour reprendre l'analogie des musiciens : la convergence uniforme, c'est lorsque, à partir d'un certain moment, l'ensemble de l'orchestre est globalement accordé avec une précision donnée. Il n'y a pas un musicien plus en retard qu'un autre dans le processus global d'accordement. L'écart maximal de justesse de l'ensemble de l'orchestre diminue de manière coordonnée.

La distinction est subtile mais capitale. La convergence simple nous assure que les choses s'améliorent localement, mais elle ne garantit rien sur la "qualité" globale de l'approximation à un instant donné. Elle peut entraîner des surprises désagréables, comme la perte de propriétés fondamentales : une suite de fonctions "lisses" (continues, par exemple) pourrait converger simplement vers une fonction qui ne l'est pas du tout. La convergence uniforme, en revanche, est le garde-fou qui préserve la "qualité" des fonctions dans le processus de passage à la limite. Elle est le prix à payer pour que les belles propriétés des fonctions initiales (continuité, intégrabilité) se transmettent à la fonction limite.

La nécessité de formaliser rigoureusement ces concepts est née d'une crise historique en analyse au XIXe siècle. Avant les travaux de figures éminentes telles qu'Augustin-Louis Cauchy et Karl Weierstrass, les mathématiciens manipulaient les séries de fonctions et les limites avec une intuition souvent géométrique ou physique, pensant qu'une limite infinie de fonctions continues conserverait naturellement la continuité. Cauchy lui-même, en 1821, affirmait à tort qu'une série de fonctions continues convergente produisait nécessairement une somme continue, ce qui fut démenti par des contre-exemples de Fourier et d'Abel. C'est Weierstrass et Gudermann qui introduiront la distinction cruciale de la "convergence uniforme" pour sauver l'édifice de l'analyse, prouvant que cette rigueur globale était la clé pour préserver la continuité, l'intégrabilité, et la dérivabilité lors d'un passage à la limite.

\section*{2. Formalisation : Le Protocole d'Exégèse Conceptuelle}

Nous allons maintenant appliquer le protocole d'exégèse conceptuelle pour formaliser ces deux types de convergence.

\subsection*{Notion 1 : La Convergence Simple}

\textbf{A. Énoncé Symbolique Strict}

Soit un ensemble $I$ (généralement un sous-ensemble de $\mathbb{R}$ ou $\mathbb{C}$, tel qu'un intervalle). Soit $(f_n)_{n \in \mathbb{N}}$ une suite de fonctions définies sur $I$ à valeurs dans $\mathbb{R}$ ou $\mathbb{C}$. Soit $f$ une fonction définie sur $I$.
On dit que la suite de fonctions $(f_n)_{n \in \mathbb{N}}$ converge simplement vers $f$ sur $I$ si :
$$ \forall x \in I, \forall \epsilon > 0, \exists N \in \mathbb{N}, \text{ tel que } \forall n \ge N, |f_n(x) - f(x)| < \epsilon $$

\textbf{B. Anatomie et Typage Chirurgical}

*   \textbf{L'ensemble de départ $I$ :} C'est le domaine de définition commun à toutes les fonctions $f_n$ et à la fonction limite $f$. Il s'agit d'une donnée globale fixe du problème.
*   \textbf{Les fonctions $f_n : I \to \mathbb{K}$ :} Pour chaque entier $n \in \mathbb{N}$, $f_n$ associe à un point $x \in I$ un scalaire $f_n(x)$ (réel ou complexe).
*   \textbf{La fonction limite $f : I \to \mathbb{K}$ :} C'est la fonction "cible".
\textit{   \textbf{L'ordre des quantificateurs (Crucial !) :} Observez attentivement le début de la proposition : $\forall x \in I, \forall \epsilon > 0, \exists N \in \mathbb{N}$. Le rang $N$ à partir duquel l'écart est inférieur à $\epsilon$ dépend }à la fois\textit{ de la précision $\epsilon$ voulue }et* du point $x$ choisi. On écrit formellement $N = N(\epsilon, x)$. C'est le marqueur absolu de la convergence simple : on étudie la convergence point par point de manière isolée.

\textbf{C. Exemples de Validation}

\textit{Exemple Trivial :}
Soit $f_n(x) = \frac{x}{n}$ sur $I = \mathbb{R}$. Pour tout $x \in \mathbb{R}$ fixé, $\lim_{n \to \infty} \frac{x}{n} = 0$. Donc la suite $(f_n)$ converge simplement vers la fonction nulle $f(x) = 0$ sur $\mathbb{R}$.
Posons le cadre formel : Soit $x \in \mathbb{R}$ et $\epsilon > 0$. On cherche $N$ tel que pour $n \ge N$, $|\frac{x}{n} - 0| < \epsilon$.
$$ \frac{|x|}{n} < \epsilon \iff n > \frac{|x|}{\epsilon} $$
Il suffit de prendre l'entier $N = \lfloor \frac{|x|}{\epsilon} \rfloor + 1$. Ce $N$ dépend explicitement de $x$ et de $\epsilon$.

\textit{Exemple Complexe :}
Considérons $f_n(x) = x^n e^{-nx}$ sur $I = [0, +\infty[$. Pour $x=0$, $f_n(0) = 0 \to 0$. Pour $x 0$, l'exponentielle l'emporte sur la puissance, donc $\lim_{n \to \infty} x^n e^{-nx} = 0$. La suite converge simplement vers la fonction nulle sur $[0, +\infty[$.

\textbf{D. Cas Pathologiques et Contre-exemples}

\textit{Contre-exemple aux limites (La bosse glissante) :}
Soit $f_n(x)$ définie sur $\mathbb{R}$ par $f_n(x) = 1$ si $x \in [n, n+1]$ et $0$ ailleurs.
Pour tout $x \in \mathbb{R}$ fixé, il existe un rang $N$ (tel que $N x$) tel que pour tout $n \ge N$, $x \notin [n, n+1]$, et donc $f_n(x) = 0$.
Ainsi, $\lim_{n \to \infty} f_n(x) = 0$ pour tout $x$. La suite converge simplement vers la fonction nulle. Pourtant, la "bosse" de hauteur 1 existe toujours, elle glisse simplement vers l'infini, montrant que la convergence simple "perd de vue" l'allure globale de la fonction.

\subsection*{Notion 2 : La Convergence Uniforme}

\textbf{A. Énoncé Symbolique Strict}

On dit que la suite de fonctions $(f_n)_{n \in \mathbb{N}}$ converge uniformément vers $f$ sur $I$ si :
$$ \forall \epsilon > 0, \exists N \in \mathbb{N}, \text{ tel que } \forall n \ge N, \forall x \in I, |f_n(x) - f(x)| < \epsilon $$
Alternativement, en introduisant la norme infini $\|g\|_{\infty, I} = \sup_{x \in I} |g(x)|$, la convergence uniforme équivaut à :
$$ \lim_{n \to \infty} \|f_n - f\|_{\infty, I} = 0 $$

\textbf{B. Anatomie et Typage Chirurgical}

\textit{   \textbf{Le basculement quantificatif (Le cœur de la différence) :} L'énoncé commence par $\forall \epsilon > 0, \exists N \in \mathbb{N}, \forall n \ge N, \forall x \in I$. Le quantificateur $\forall x \in I$ a été repoussé }après\textit{ l'existence de $N$. Cela signifie que le rang $N$ dépend de $\epsilon$, mais \textbf{ne dépend plus de $x$}. $N = N(\epsilon)$. C'est un rang global, universel pour tout l'intervalle $I$. Dès que $n \ge N$, l'écart $|f_n(x) - f(x)|$ est inférieur à $\epsilon$ }simultanément* pour absolument tous les points de $I$.
*   \textbf{La notion de bande de tolérance (Le "Tube") :} Dire que $\forall x \in I, |f_n(x) - f(x)| < \epsilon$ revient à affirmer que le graphe entier de la fonction $f_n$ est entièrement contenu dans un "tube" de largeur $2\epsilon$ centré sur le graphe de la fonction limite $f$ : $f(x) - \epsilon < f_n(x) < f(x) + \epsilon$.
*   \textbf{La norme sup $\| \cdot \|_{\infty, I}$ :} Cette norme (le supremum des valeurs absolues sur $I$) mesure précisément le pire écart entre $f_n$ et $f$ sur tout l'intervalle. La convergence uniforme stipule que ce pire écart absolu tend vers zéro lorsque $n$ tend vers l'infini.

\textbf{C. Exemples de Validation}

\textit{Exemple Trivial :}
Soit $f_n(x) = \frac{\sin(nx)}{n}$ sur $I = \mathbb{R}$. La fonction limite simple est clairement $f(x) = 0$.
Pour la convergence uniforme, cherchons le supremum de l'écart :
$|f_n(x) - f(x)| = \left| \frac{\sin(nx)}{n} \right| \le \frac{1}{n}$ pour tout $x \in \mathbb{R}$.
Donc, $\|f_n - f\|_{\infty, \mathbb{R}} \le \frac{1}{n}$.
Puisque $\lim_{n \to \infty} \frac{1}{n} = 0$, on a $\lim_{n \to \infty} \|f_n - f\|_{\infty, \mathbb{R}} = 0$. La convergence est bien uniforme sur $\mathbb{R}$.

\textit{Exemple de Non-Uniformité (La cassure de la continuité) :}
Soit $f_n(x) = x^n$ sur $I = [0, 1]$.
Convergence simple (déjà vue dans l'introduction) :
- Pour $x \in [0, 1[$, $x^n \to 0$.
- Pour $x = 1$, $1^n = 1 \to 1$.
La fonction limite $f$ vaut $0$ sur $[0, 1[$ et $1$ en $1$. $f$ est discontinue en $1$.
Étudions la convergence uniforme. Les fonctions $f_n$ sont continues sur $[0,1]$. La fonction limite $f$ ne l'est pas. Le théorème de continuité stipule que si une suite de fonctions continues converge uniformément, la limite DOIT être continue. Par contraposée, puisque la limite est discontinue, la convergence \textbf{ne peut pas être uniforme}.
Vérifions formellement par la norme sup :
Sur l'intervalle $[0, 1[$, l'écart est $|f_n(x) - f(x)| = |x^n - 0| = x^n$.
$\|f_n - f\|_{\infty, [0, 1[} = \sup_{x \in [0, 1[} x^n = 1$.
La limite de ce supremum lorsque $n \to \infty$ est $1$, et non $0$. La convergence n'est pas uniforme.

\textbf{D. Cas Pathologiques et Contre-exemples}

\textit{Contre-exemple de la limite continue sans convergence uniforme :}
Soit $f_n(x) = nx(1-x)^n$ sur $I = [0, 1]$.
Pour $x=0$, $f_n(0) = 0 \to 0$.
Pour $x \in ]0, 1]$, on reconnait la forme $n q^n$ avec $|q| = 1-x < 1$. D'après les croissances comparées, $\lim_{n \to \infty} n(1-x)^n = 0$.
La suite converge donc simplement vers la fonction nulle $f(x) = 0$ sur $[0, 1]$. La fonction limite est continue. La continuité est préservée.
Cependant, la convergence est-elle uniforme ? Calculons le supremum.
La dérivée $f_n'(x) = n(1-x)^n - n^2x(1-x)^{n-1} = n(1-x)^{n-1}(1 - x - nx) = n(1-x)^{n-1}(1 - (n+1)x)$.
La dérivée s'annule en $x_n = \frac{1}{n+1}$, qui est le point de maximum sur $[0, 1]$.
La valeur maximale est $f_n\left(\frac{1}{n+1}\right) = n \left(\frac{1}{n+1}\right) \left(1 - \frac{1}{n+1}\right)^n = \frac{n}{n+1} \left(\frac{n}{n+1}\right)^n = \left(\frac{n}{n+1}\right)^{n+1}$.
$\|f_n - f\|_{\infty, [0,1]} = \left(1 - \frac{1}{n+1}\right)^{n+1}$.
Or, $\lim_{n \to \infty} \left(1 - \frac{1}{n+1}\right)^{n+1} = e^{-1} \approx 0.368 \neq 0$.
Le supremum de l'erreur ne tend pas vers zéro. Une "bosse" de hauteur $\frac{1}{e}$ se concentre de plus en plus près de 0 mais ne disparaît jamais. La convergence est simple, la limite est continue, \textbf{mais la convergence n'est pas uniforme}. C'est le cas pathologique par excellence qui montre que la convergence uniforme est une condition forte.

\section*{3. Démonstrations Pas-à-Pas}

\subsection*{Démonstration du Théorème Pivot : Le Théorème d'Interversion des Limites (Continuité de la Limite)}

\textbf{Théorème de Continuité de la Fonction Limite :}
Soit $I$ un intervalle de $\mathbb{R}$. Soit $(f_n)_{n \in \mathbb{N}}$ une suite de fonctions définies sur $I$, convergeant uniformément vers une fonction $f$ sur $I$.
Si, pour tout entier $n \in \mathbb{N}$, la fonction $f_n$ est continue en un point $a \in I$, alors la fonction limite $f$ est également continue en ce point $a$.

1. \textbf{Initialisation / Cadre :}
   L'objectif est de démontrer la continuité de la fonction limite $f$ au point $a \in I$.
   Rappel de l'objectif formel de continuité en $a$ :
   On doit prouver que : $\forall \epsilon > 0, \exists \delta > 0, \forall x \in I, |x - a| < \delta \implies |f(x) - f(a)| < \epsilon$.
   Nous disposons de deux hypothèses fondamentales :
   - Hypothèse 1 (Convergence Uniforme) : $\forall \epsilon' 0, \exists N \in \mathbb{N}, \forall n \ge N, \forall x \in I, |f_n(x) - f(x)| < \epsilon'$.
   - Hypothèse 2 (Continuité des $f_n$) : Pour un $n$ fixé, $\forall \epsilon'' 0, \exists \delta > 0, \forall x \in I, |x - a| < \delta \implies |f_n(x) - f_n(a)| < \epsilon''$.

   La stratégie consiste à majorer l'écart cible $|f(x) - f(a)|$ en introduisant astucieusement les fonctions $f_n$ par une "astuce des trois tiers" basée sur l'inégalité triangulaire.

2. \textbf{Étape 1 : Décomposition par l'Inégalité Triangulaire}
   Pour tout $x \in I$ et tout entier $n \in \mathbb{N}$, nous pouvons écrire l'identité tautologique suivante en ajoutant et soustrayant $f_n(x)$ et $f_n(a)$ à l'intérieur de la valeur absolue :
   $$ f(x) - f(a) = f(x) - f_n(x) + f_n(x) - f_n(a) + f_n(a) - f(a) $$
   En appliquant l'inégalité triangulaire ($|A + B + C| \le |A| + |B| + |C|$), nous obtenons la majoration fondamentale :
   $$ |f(x) - f(a)| \le |f(x) - f_n(x)| + |f_n(x) - f_n(a)| + |f_n(a) - f(a)| $$
   Cette ligne est la pierre angulaire de la preuve. Elle décompose l'erreur totale en trois composantes que nous allons contrôler séparément.

3. \textbf{Étape 2 (Transition micro-calculatoire) : Majoration des trois composantes}
   Soit un scalaire de précision arbitraire $\epsilon > 0$. Nous allons allouer un "budget d'erreur" de $\frac{\epsilon}{3}$ à chacune des trois composantes de la somme.

   \textbf{Contrôle du premier et du troisième terme (Convergence Uniforme) :}
   D'après l'Hypothèse 1 (Convergence Uniforme de $(f_n)$ vers $f$), appliquée avec le paramètre de tolérance $\epsilon' = \frac{\epsilon}{3}$ :
   Il existe un entier $N \in \mathbb{N}$ tel que, pour tout entier $n \ge N$ et pour \textit{tout point} $t \in I$ (donc a fortiori pour $t=x$ et pour $t=a$) :
   $$ |f_n(t) - f(t)| < \frac{\epsilon}{3} $$
   Fixons définitivement un entier $n_0 \ge N$ (par exemple $n_0 = N$). L'inégalité précédente est vraie pour la fonction spécifique $f_{n_0}$ :
   - Pour le premier terme : $|f(x) - f_{n_0}(x)| = |f_{n_0}(x) - f(x)| < \frac{\epsilon}{3}$, et ce, quel que soit $x \in I$.
   - Pour le troisième terme : $|f_{n_0}(a) - f(a)| < \frac{\epsilon}{3}$.

   À ce stade de la preuve, l'entier $n_0$ est \textbf{fixé}.

   \textbf{Contrôle du deuxième terme (Continuité de $f_{n_0}$) :}
   Considérons la fonction $f_{n_0}$. D'après l'Hypothèse 2 (Continuité des $f_n$), la fonction $f_{n_0}$ est continue au point $a$.
   En appliquant la définition de la continuité avec le paramètre de tolérance $\epsilon'' = \frac{\epsilon}{3}$, nous savons qu'il existe un rayon $\delta > 0$ tel que :
   $$ \forall x \in I, \text{ si } |x - a| < \delta \text{ alors } |f_{n_0}(x) - f_{n_0}(a)| < \frac{\epsilon}{3} $$

4. \textbf{Conclusion : Assemblage final}
   Synthétisons les résultats des étapes précédentes.
   Soit $\epsilon > 0$. Nous avons construit un rayon $\delta > 0$ tel que, pour tout $x \in I$ vérifiant $|x - a| < \delta$, les trois majorations suivantes sont simultanément vraies :
   1. $|f(x) - f_{n_0}(x)| < \frac{\epsilon}{3}$
   2. $|f_{n_0}(x) - f_{n_0}(a)| < \frac{\epsilon}{3}$
   3. $|f_{n_0}(a) - f(a)| < \frac{\epsilon}{3}$

   En réinjectant ces bornes dans l'inégalité triangulaire de l'Étape 1 :
   $$ |f(x) - f(a)| \le |f(x) - f_{n_0}(x)| + |f_{n_0}(x) - f_{n_0}(a)| + |f_{n_0}(a) - f(a)| $$
   $$ |f(x) - f(a)| < \frac{\epsilon}{3} + \frac{\epsilon}{3} + \frac{\epsilon}{3} $$
   $$ |f(x) - f(a)| < \epsilon $$

   Nous avons rigoureusement établi que : $\forall \epsilon > 0, \exists \delta > 0, \forall x \in I, |x - a| < \delta \implies |f(x) - f(a)| < \epsilon$.
   Ceci est l'exacte définition de la continuité de la fonction $f$ au point $a$. Le théorème est prouvé.
   \textbf{C.Q.F.D.}

\section*{4. Exercices d'Application}

[Les exercices détaillés se trouvent dans les fichiers exos/Exo-01.md à exos/Exo-10.md, qui appliquent tous la rigueur calculatoire pas-à-pas attendue].

\section*{5. Application en Intelligence Artificielle}

Les concepts de convergence des suites de fonctions jouent un rôle discret mais fondamental dans le développement et l'analyse des algorithmes d'Intelligence Artificielle, en particulier dans l'apprentissage automatique.

1.  \textbf{Apprentissage des Réseaux de Neurones :} Un réseau de neurones est essentiellement une fonction paramétrée, $h_{\theta}(x)$, où $\theta$ représente l'ensemble des poids et biais. Le processus d'entraînement consiste à ajuster ces paramètres $\theta$ de manière itérative (par exemple, via la descente de gradient stochastique) pour minimiser une fonction de coût. Chaque itération de l'entraînement produit un nouvel ensemble de paramètres $\theta_k$, et donc une nouvelle fonction $h_{\theta_k}(x)$. La "convergence" d'un modèle d'IA peut être interprétée comme la convergence de cette suite de fonctions $(h_{\theta_k}(x))$ vers une fonction optimale $h_{\theta^*}(x)$.
    \textit{   \textbf{Convergence simple (ou ponctuelle) :} Si, après un certain nombre d'époques, le réseau fournit des prédictions de plus en plus précises pour }chaque* exemple d'entraînement (ou de test) individuel, on peut parler de convergence simple. Cela garantit que l'erreur sur un point de donnée spécifique tend vers zéro.
    \textit{   \textbf{Convergence uniforme :} Ce qui est souvent désiré en IA, c'est que l'erreur du réseau diminue }uniformément* sur l'ensemble du jeu de données (entraînement, validation ou test). Cela signifie que le réseau ne s'améliore pas seulement sur certains exemples au détriment d'autres, mais que sa performance globale s'améliore, et que l'écart maximal entre ses prédictions et les vraies valeurs sur l'ensemble du jeu de données tend vers zéro. Des techniques comme la régularisation, le "batch normalization" ou le "dropout" visent implicitement à favoriser une convergence plus uniforme et une meilleure généralisation en évitant l'hyper-spécialisation sur des sous-ensembles de données.

2.  \textbf{Théorèmes d'Approximation Universelle :} Ces théorèmes sont fondamentaux pour justifier la capacité des réseaux de neurones à modéliser des fonctions complexes. Le théorème d'approximation universelle de Cybenko ou de Hornik stipule qu'un réseau de neurones feedforward avec une seule couche cachée et un nombre suffisant de neurones peut approximer n'importe quelle fonction continue sur un compact avec une précision arbitraire. Cette "précision arbitraire" est une formulation directe de la convergence uniforme : pour toute marge d'erreur $\epsilon > 0$, il existe un réseau (une fonction) qui s'approche uniformément de la fonction cible à moins de $\epsilon$.

\section*{6. Liens Sémantiques}

- \textbf{Concepts Précédents requis :} Jalon 14 (Suites réelles et complexes), Jalon 18 (Continuité des fonctions d'une variable réelle)
- \textbf{Concepts Futurs dépendants :} Jalon 22 (Séries de fonctions), Jalon 23 (Séries entières)

\newpage
\chapter*{Exercices d\'Application}
\addcontentsline{toc}{chapter}{Exercices d'Application}


\chapter*{Exercice 1 : Convergence simple de $x^n$}

\textbf{Niveau :} $\star \circ \circ \circ \circ $

\section*{Problème}

Étudier la convergence simple de $f_n(x) = x^n$ sur $[0,1]$.

\section*{Démonstration et Solution}

Pour étudier la convergence simple de la suite de fonctions $f_n(x) = x^n$ sur le segment $[0,1]$, nous devons fixer un réel $x \in [0,1]$ arbitraire et étudier la limite de la suite numérique $(f_n(x))_{n \in \mathbb{N}}$.

Nous distinguons deux cas rigoureux :

\textbf{Cas 1 : $x \in [0, 1[$}
Si $x \in [0, 1[$, alors la valeur de $x$ est strictement inférieure à 1. La suite numérique $(x^n)_{n \in \mathbb{N}}$ est une suite géométrique de raison $q = x$. Puisque $0 \leq x < 1$, nous avons $|q| < 1$. D'après le théorème sur les limites des suites géométriques, si la raison appartient à l'intervalle ouvert $]-1, 1[$, la limite de $q^n$ lorsque $n$ tend vers l'infini est strictement égale à 0.
Ainsi, pour tout $x \in [0, 1[$, $\lim_{n \to \infty} f_n(x) = \lim_{n \to \infty} x^n = 0$.

\textbf{Cas 2 : $x = 1$}
Si $x = 1$, alors pour tout entier naturel $n$, $f_n(1) = 1^n = 1$. La suite numérique considérée est donc la suite constante égale à 1. La limite d'une suite constante est sa valeur, donc :
$\lim_{n \to \infty} f_n(1) = 1$.

\textbf{Conclusion de l'étude de la convergence simple :}
La suite de fonctions $(f_n)$ converge simplement sur $[0,1]$ vers une fonction limite $f$ définie par :
$f(x) = 0$ si $x \in [0, 1[$
$f(x) = 1$ si $x = 1$

\newpage


\chapter*{Exercice 2 : Convergence simple de x/n}

\textbf{Niveau :} $\star \circ \circ \circ \circ $

\section*{Problème}

Étudier la convergence simple de $f_n(x) = x/n$ sur $\mathbb{R}$.

\section*{Démonstration et Solution}

Pour étudier la convergence simple de $f_n(x) = \frac{x}{n}$ sur $\mathbb{R}$, fixons un réel arbitraire $x \in \mathbb{R}$.
Nous considérons la suite numérique définie par $u_n = \frac{x}{n}$ pour $n \geq 1$.

Par définition de la limite, $\lim_{n \to \infty} \frac{1}{n} = 0$.
Le réel $x$ étant fixé, il s'agit d'une constante par rapport à la variable d'indice $n$. Par les propriétés des opérations sur les limites (produit d'une limite finie par une constante), nous obtenons :
$\lim_{n \to \infty} u_n = x \times \lim_{n \to \infty} \frac{1}{n} = x \times 0 = 0$.

Ce résultat est vrai quelle que soit la valeur initiale fixée pour le réel $x$.
Ainsi, pour tout $x \in \mathbb{R}$, la suite numérique $(f_n(x))$ converge vers 0.

\textbf{Conclusion :}
La suite de fonctions $(f_n)$ converge simplement sur $\mathbb{R}$ vers la fonction limite nulle, c'est-à-dire la fonction $f$ définie pour tout $x \in \mathbb{R}$ par $f(x) = 0$.

\newpage


\chapter*{Exercice 3 : Convergence uniforme de x/n}

\textbf{Niveau :} $\star \star \circ \circ \circ $

\section*{Problème}

Montrer que $f_n(x) = x/n$ ne converge pas uniformément sur $\mathbb{R}$.

\section*{Démonstration et Solution}

D'après l'exercice précédent, la suite de fonctions $f_n(x) = \frac{x}{n}$ converge simplement vers la fonction nulle $f(x) = 0$ sur $\mathbb{R}$.
Pour déterminer si cette convergence est uniforme sur $\mathbb{R}$, nous devons évaluer la norme de la convergence uniforme, c'est-à-dire le supremum de l'écart absolu entre les fonctions $f_n$ et la fonction limite $f$ sur l'ensemble du domaine de définition.

Posons la norme infinie de la différence :
$\|f_n - f\|_\infty = \sup_{x \in \mathbb{R}} |f_n(x) - f(x)|$
En remplaçant par nos fonctions :
$\|f_n - f\|_\infty = \sup_{x \in \mathbb{R}} \left| \frac{x}{n} - 0 \right| = \sup_{x \in \mathbb{R}} \frac{|x|}{n}$

Pour un entier $n \geq 1$ fixé, regardons le comportement de la fonction $x \mapsto \frac{|x|}{n}$ sur $\mathbb{R}$.
Lorsque $x$ tend vers $+\infty$, la quantité $\frac{|x|}{n}$ tend vers $+\infty$.
L'ensemble des valeurs $\left\{ \frac{|x|}{n} \mid x \in \mathbb{R} \right\}$ n'est donc pas majoré. Son supremum dans $\mathbb{R} \cup \{+\infty\}$ est donc infini.
Par conséquent, pour tout $n \geq 1$, $\|f_n - f\|_\infty = +\infty$.

Puisque la suite des normes $\|f_n - f\|_\infty$ ne tend pas vers 0 (elle est constante égale à $+\infty$), la convergence n'est par définition pas uniforme sur l'ensemble $\mathbb{R}$.

\newpage


\chapter*{Exercice 4 : Convergence uniforme sur segment}

\textbf{Niveau :} $\star \star \circ \circ \circ $

\section*{Problème}

Montrer que $f_n(x) = x/n$ converge uniformément sur tout segment $[-M, M]$ avec $M > 0$.

\section*{Démonstration et Solution}

Considérons le segment fini $[-M, M]$ où $M$ est un réel strictement positif.
Nous savons que la suite $f_n(x) = \frac{x}{n}$ converge simplement vers la fonction nulle $f(x) = 0$ sur ce segment (puisque cela est vrai sur tout $\mathbb{R}$).

Évaluons maintenant la norme de la différence sur ce segment restreint :
$\|f_n - f\|_{\infty, [-M,M]} = \sup_{x \in [-M, M]} |f_n(x) - f(x)| = \sup_{x \in [-M, M]} \frac{|x|}{n}$

Pour tout $x \in [-M, M]$, nous avons l'inégalité stricte ou large : $|x| \leq M$.
Par conséquent, en divisant par l'entier strictement positif $n$, nous obtenons pour tout $x \in [-M, M]$ :
$\frac{|x|}{n} \leq \frac{M}{n}$
Puisque la borne $M/n$ est atteinte en posant $x=M$ (ou $x=-M$), le supremum est exactement égal à cette valeur :
$\|f_n - f\|_{\infty, [-M,M]} = \frac{M}{n}$

Étudions la limite de cette suite de normes :
$\lim_{n \to \infty} \|f_n - f\|_{\infty, [-M,M]} = \lim_{n \to \infty} \frac{M}{n} = M \times \lim_{n \to \infty} \frac{1}{n} = M \times 0 = 0$.

\textbf{Conclusion :}
La norme infinie de la différence entre $f_n$ et $f$ sur le segment $[-M, M]$ tend rigoureusement vers 0 lorsque $n$ tend vers l'infini. Par définition, cela signifie que la suite de fonctions $(f_n)$ converge uniformément vers la fonction nulle sur le segment $[-M, M]$.

\newpage


\chapter*{Exercice 5 : Fonction limite discontinue}

\textbf{Niveau :} $\star \star \star \circ \circ $

\section*{Problème}

Étudier la convergence de $f_n(x) = x^n$ sur $[0,1]$. La fonction limite est-elle continue ? Que peut-on dire de la convergence uniforme ?

\section*{Démonstration et Solution}

\textbf{Étape 1 : Étude de la convergence simple}
Nous avons démontré à l'Exercice 1 que la suite de fonctions $f_n(x) = x^n$ converge simplement sur $[0,1]$ vers la fonction limite $f$ définie par :
$f(x) = 0$ si $x \in [0, 1[$
$f(x) = 1$ si $x = 1$.

\textbf{Étape 2 : Continuité de la fonction limite}
Pour tout entier $n \geq 1$, la fonction $f_n(x) = x^n$ est une fonction polynôme, elle est donc parfaitement continue sur l'intervalle compact $[0,1]$.
Cependant, analysons la continuité de la fonction limite $f$ au point $x_0 = 1$.
Calculons la limite à gauche de $f$ en 1 :
Pour tout $x \in [0, 1[$, $f(x) = 0$. Donc, $\lim_{x \to 1^-} f(x) = \lim_{x \to 1^-} 0 = 0$.
Or, la valeur de la fonction limite en 1 est $f(1) = 1$.
Puisque $\lim_{x \to 1^-} f(x) \neq f(1)$, la fonction $f$ présente une discontinuité de première espèce au point $x=1$. Elle n'est donc pas continue sur $[0,1]$.

\textbf{Étape 3 : Conséquence sur la convergence uniforme}
Le théorème fondamental de la continuité des limites de suites de fonctions stipule que si une suite de fonctions continues $(f_n)$ converge uniformément vers une fonction $f$ sur un intervalle $I$, alors la fonction limite $f$ est obligatoirement continue sur cet intervalle $I$.
Nous raisonnons ici par contraposée :
Puisque toutes les fonctions de la suite $f_n$ sont continues sur $[0,1]$, et puisque la fonction limite $f$ \textbf{n'est pas} continue sur $[0,1]$, la convergence de la suite $(f_n)$ vers $f$ \textbf{ne peut absolument pas} être uniforme sur $[0,1]$.

\newpage


\chapter*{Exercice 6 : Théorème d'interversion limite-intégrale (Contre-exemple)}

\textbf{Niveau :} $\star \star \star \circ \circ $

\section*{Problème}

Considérer $f_n(x) = n^2 x (1-x^2)^n$ sur $[0,1]$. Étudier la convergence simple. Comparer l'intégrale de la limite et la limite de l'intégrale.

\section*{Démonstration et Solution}

\textbf{Étape 1 : Étude de la convergence simple}
Soit $f_n(x) = n^2 x (1-x^2)^n$ sur $[0,1]$.
Si $x = 0$, $f_n(0) = n^2 \times 0 \times 1^n = 0$. La limite est 0.
Si $x = 1$, $f_n(1) = n^2 \times 1 \times 0^n = 0$. La limite est 0.
Si $x \in ]0, 1[$, alors $1-x^2 \in ]0, 1[$. Posons $q = 1-x^2$, avec $0 < q < 1$.
L'expression devient $f_n(x) = x \times n^2 q^n$.
Puisque $0 < q < 1$, nous avons une forme indéterminée du type "+$\infty \times 0$". Par le théorème de croissance comparée entre les fonctions puissances et les exponentielles (ou en écrivant $q^n = e^{n \ln q}$ avec $\ln q < 0$), l'exponentielle l'emporte. Rigoureusement, $\lim_{n \to \infty} n^k q^n = 0$ pour tout $k > 0$ et $q \in ]0,1[$.
Ainsi, $\lim_{n \to \infty} f_n(x) = x \times 0 = 0$.
La suite converge simplement vers la fonction nulle $f(x)=0$ sur $[0,1]$.

\textbf{Étape 2 : Intégrale de la limite}
La fonction limite étant nulle partout, son intégrale sur le segment est nulle :
$\int_0^1 f(x) dx = \int_0^1 0 dx = 0$.

\textbf{Étape 3 : Limite de l'intégrale}
Calculons l'intégrale de $f_n(x)$ sur $[0,1]$ pour un $n$ fixé :
$I_n = \int_0^1 n^2 x (1-x^2)^n dx$
Effectuons le changement de variable $u = 1-x^2$.
Alors $du = -2x dx$, ce qui donne $x dx = -\frac{1}{2} du$.
Les bornes deviennent : pour $x=0$, $u=1$ ; pour $x=1$, $u=0$.
$I_n = \int_1^0 n^2 (u)^n \left(-\frac{1}{2}\right) du = \frac{n^2}{2} \int_0^1 u^n du = \frac{n^2}{2} \left[ \frac{u^{n+1}}{n+1} \right]_0^1 = \frac{n^2}{2(n+1)}$
Calculons la limite de $I_n$ lorsque $n$ tend vers l'infini :
$\lim_{n \to \infty} I_n = \lim_{n \to \infty} \frac{n^2}{2n+2} = \lim_{n \to \infty} \frac{n^2}{2n(1 + 1/n)} = \lim_{n \to \infty} \frac{n}{2(1+1/n)} = +\infty$.

\textbf{Conclusion :}
Nous observons que la limite des intégrales est $+\infty$, tandis que l'intégrale de la limite est $0$. Les deux quantités sont distinctes ($\lim \int f_n \neq \int \lim f_n$). Ceci prouve formellement, par l'absurde via le théorème d'interversion, que la convergence de $(f_n)$ vers $f$ n'est pas uniforme sur $[0,1]$.

\newpage


\chapter*{Exercice 7 : Convergence uniforme de arctan(nx)}

\textbf{Niveau :} $\star \star \star \star \circ $

\section*{Problème}

Étudier la convergence de $f_n(x) = \frac{1}{n}\arctan(nx)$ sur $\mathbb{R}$.

\section*{Démonstration et Solution}

Considérons la suite de fonctions $f_n(x) = \frac{1}{n} \arctan(nx)$ sur $\mathbb{R}$.

\textbf{Étape 1 : Convergence simple}
Pour tout réel $x \in \mathbb{R}$, fixons $x$ et étudions la limite quand $n \to \infty$.
Nous connaissons les bornes de la fonction arc tangente : pour tout réel $y$, $|\arctan(y)| < \frac{\pi}{2}$.
Ainsi, pour tout réel $x$ et tout entier $n \geq 1$ :
$|f_n(x)| = \left| \frac{1}{n} \arctan(nx) \right| = \frac{1}{n} |\arctan(nx)| < \frac{1}{n} \times \frac{\pi}{2}$
Puisque $\lim_{n \to \infty} \frac{\pi}{2n} = 0$, le théorème des gendarmes implique que $\lim_{n \to \infty} f_n(x) = 0$.
La suite converge simplement vers la fonction nulle $f(x)=0$ sur $\mathbb{R}$.

\textbf{Étape 2 : Convergence uniforme}
Pour étudier la convergence uniforme, déterminons la norme infinie de la différence $\|f_n - f\|_\infty$ sur $\mathbb{R}$.
$\|f_n - f\|_\infty = \sup_{x \in \mathbb{R}} |f_n(x) - 0| = \sup_{x \in \mathbb{R}} \frac{|\arctan(nx)|}{n}$
Nous savons que la fonction $x \mapsto \arctan(nx)$ est strictement croissante et impaire. Son supremum sur $\mathbb{R}$ est atteint asymptotiquement lorsque $x \to +\infty$.
$\sup_{x \in \mathbb{R}} |\arctan(nx)| = \lim_{x \to +\infty} \arctan(nx) = \frac{\pi}{2}$
Par conséquent, en divisant par la constante positive $n$, nous obtenons la valeur exacte de la norme infinie :
$\|f_n - f\|_\infty = \frac{1}{n} \times \frac{\pi}{2} = \frac{\pi}{2n}$
Étudions maintenant la limite de cette norme lorsque $n$ tend vers l'infini :
$\lim_{n \to \infty} \|f_n - f\|_\infty = \lim_{n \to \infty} \frac{\pi}{2n} = 0$.
Puisque la norme infinie de l'écart tend vers 0, la convergence est, par définition, uniforme sur l'ensemble $\mathbb{R}$.

\newpage


\chapter*{Exercice 8 : Suite de fonctions continues vers continue sans CU}

\textbf{Niveau :} $\star \star \star \star \circ $

\section*{Problème}

Construire une suite de fonctions continues sur $[0,1]$ qui converge simplement vers $0$, mais dont la norme infinie vaut toujours $1$.

\section*{Démonstration et Solution}

Pour répondre à la question, nous devons concevoir une suite de fonctions "chapeau" ou "pointe de tente" dont la hauteur reste constante égale à 1, mais dont la base se rétrécit vers 0 à mesure que $n$ augmente.

Définissons pour chaque entier $n \geq 2$ la fonction $f_n$ sur $[0,1]$ par des segments de droite reliant les points suivants :
- Point $A : (0, 0)$
- Point $B : (\frac{1}{2n}, 1)$
- Point $C : (\frac{1}{n}, 0)$
- Point $D : (1, 0)$
Explicitement, cela s'écrit formellement par morceaux :
$f_n(x) = 2nx$ si $x \in [0, \frac{1}{2n}[$
$f_n(x) = -2nx + 2$ si $x \in [\frac{1}{2n}, \frac{1}{n}[$
$f_n(x) = 0$ si $x \in [\frac{1}{n}, 1]$.
Par construction (raccordement continu aux points de rupture), chaque fonction $f_n$ est continue sur le segment $[0,1]$.

\textbf{Étape 1 : Convergence simple}
Fixons $x \in [0,1]$.
Si $x = 0$, $f_n(0) = 0$ pour tout $n$, donc la limite est 0.
Si $x > 0$, alors selon la propriété d'Archimède, il existe un entier $N$ tel que $\frac{1}{N} < x$.
Par conséquent, pour tout $n \geq N$, nous avons $\frac{1}{n} \leq \frac{1}{N} < x$.
Selon la définition de notre fonction par morceaux, si $x > \frac{1}{n}$, alors $f_n(x) = 0$.
Donc, pour un $x > 0$ fixé, la suite $f_n(x)$ est stationnaire à 0 à partir du rang $N$. Sa limite est donc 0.
La suite $(f_n)$ converge simplement vers la fonction limite continue $f(x)=0$ sur $[0,1]$.

\textbf{Étape 2 : Absence de convergence uniforme}
Calculons la norme infinie $\|f_n - f\|_\infty$ sur $[0,1]$.
Puisque $f(x)=0$, cela revient à trouver le supremum de $f_n$.
Par construction, le sommet du "chapeau" est situé à l'abscisse $x = \frac{1}{2n}$, et la valeur de la fonction en ce point est exactement $1$.
Donc, $\sup_{x \in [0,1]} |f_n(x) - 0| = f_n(\frac{1}{2n}) = 1$.
Ainsi, pour tout $n \geq 2$, $\|f_n - f\|_\infty = 1$.
La limite de la norme infinie est $\lim_{n \to \infty} 1 = 1 \neq 0$.
La convergence n'est donc pas uniforme sur $[0,1]$, bien que la suite soit constituée de fonctions continues convergeant vers une fonction continue.

\newpage


\chapter*{Exercice 9 : Théorème de Dini}

\textbf{Niveau :} $\star \star \star \star \star $

\section*{Problème}

Démontrer le théorème de Dini : si une suite $(f_n)$ de fonctions continues réelles converge simplement sur un compact $K$ vers une fonction continue $f$, et si la suite $(f_n)$ est monotone croissante (resp. décroissante), alors la convergence est uniforme.

\section*{Démonstration et Solution}

\textbf{Théorème de Dini :} Soit $K$ un espace compact (par exemple un segment fermé borné de $\mathbb{R}$). Soit $(f_n)_{n \in \mathbb{N}}$ une suite de fonctions continues de $K$ dans $\mathbb{R}$. Si $(f_n)$ converge simplement sur $K$ vers une fonction continue $f$, et si la suite $(f_n)$ est monotone en tout point (par exemple décroissante : $\forall x \in K, \forall n, f_{n+1}(x) \leq f_n(x)$), alors la convergence est uniforme sur $K$.

\textbf{Démonstration formelle (Cas décroissant) :}
Définissons la suite de fonctions auxiliaire $g_n = f_n - f$.
Puisque chaque $f_n$ est continue et $f$ est continue, par combinaison linéaire, chaque $g_n$ est une fonction continue sur le compact $K$.
La suite $(f_n)$ étant décroissante vers $f$, nous avons pour tout $x \in K$ et tout entier $n$, $f_n(x) \geq f_{n+1}(x) \geq f(x)$.
En soustrayant $f(x)$, cela implique $g_n(x) \geq g_{n+1}(x) \geq 0$. La suite $(g_n)$ est donc une suite de fonctions continues, positives et décroissante vers 0.
Démontrer la convergence uniforme de $(f_n)$ vers $f$ équivaut à démontrer la convergence uniforme de $(g_n)$ vers la fonction nulle 0.
Soit un réel arbitraire $\epsilon > 0$.
Pour chaque entier $n \in \mathbb{N}$, définissons l'ensemble de niveau $O_n$ :
$O_n = \{x \in K \mid g_n(x) < \epsilon\}$.
Puisque la fonction $g_n$ est continue, et que l'intervalle $]-\infty, \epsilon[$ est un sous-ensemble ouvert de $\mathbb{R}$, l'image réciproque $O_n = g_n^{-1}(]-\infty, \epsilon[)$ est un ensemble ouvert de $K$.
De plus, comme la suite $(g_n)$ est décroissante, si $g_n(x) < \epsilon$, alors a fortiori $g_{n+1}(x) \leq g_n(x) < \epsilon$. Cela implique rigoureusement l'inclusion ensembliste : $O_n \subset O_{n+1}$. La suite d'ouverts $(O_n)$ est croissante.
Montrons que ces ouverts recouvrent $K$. Pour tout $x \in K$, par la convergence simple de $g_n$ vers 0, on a $\lim_{n \to \infty} g_n(x) = 0$. Donc, par définition de la limite, il existe un rang $N_x$ (qui dépend de $x$) tel que $g_{N_x}(x) < \epsilon$, ce qui signifie que $x \in O_{N_x}$.
Ainsi, l'union infinie $\bigcup_{n \in \mathbb{N}} O_n$ couvre intégralement $K$.
La famille $(O_n)_{n \in \mathbb{N}}$ constitue donc un recouvrement ouvert du compact $K$.
Par la propriété fondamentale de Borel-Lebesgue (définition de la compacité), il est possible d'en extraire un sous-recouvrement fini. Il existe donc un ensemble fini d'indices $\{n_1, n_2, \dots, n_k\}$ tel que $K \subset O_{n_1} \cup O_{n_2} \cup \dots \cup O_{n_k}$.
Posons $N = \max(n_1, n_2, \dots, n_k)$. Puisque la suite d'ouverts est emboîtée de manière croissante ($O_p \subset O_q$ pour $p < q$), l'union finie est égale à son plus grand élément : $O_{n_1} \cup \dots \cup O_{n_k} \subset O_N$.
Nous obtenons alors $K \subset O_N$.
Cela signifie que pour tout $x \in K$, on a $x \in O_N$, c'est-à-dire rigoureusement $0 \leq g_N(x) < \epsilon$.
Enfin, par la décroissance de la suite $(g_n)$, pour tout $n \geq N$ et pour tout $x \in K$, on a $0 \leq g_n(x) \leq g_N(x) < \epsilon$.
En passant au supremum, on obtient que pour tout $n \geq N$, $\sup_{x \in K} g_n(x) \leq \epsilon$.
Ceci est la définition stricte de la limite uniforme de $(g_n)$ vers 0, prouvant le théorème.

\newpage


\chapter*{Exercice 10 : Équation différentielle et limite uniforme}

\textbf{Niveau :} $\star \star \star \star \star $

\section*{Problème}

Montrer que si $(f_n)$ est une suite de fonctions $C^1$ sur $[a,b]$ convergeant au moins en un point $x_0 \in [a,b]$, et si la suite des dérivées $(f_n')$ converge uniformément vers une fonction $g$, alors $(f_n)$ converge uniformément vers une fonction $f$ différentiable de dérivée $g$.

\section*{Démonstration et Solution}

\textbf{Théorème :} Soit $(f_n)$ une suite de fonctions de classe $C^1$ sur un segment $[a,b]$. Si la suite numérique $(f_n(x_0))$ converge pour au moins un point $x_0 \in [a,b]$, et si la suite des dérivées $(f_n')$ converge uniformément sur $[a,b]$ vers une fonction $g$, alors la suite $(f_n)$ converge uniformément sur $[a,b]$ vers une fonction $f$ qui est différentiable, et sa dérivée est exactement $f' = g$.

\textbf{Démonstration formelle :}
Puisque les fonctions $f_n$ sont de classe $C^1$, le Théorème Fondamental de l'Analyse nous permet d'écrire, pour tout point $x \in [a,b]$ et par rapport au point fixe $x_0$ :
$f_n(x) = f_n(x_0) + \int_{x_0}^x f_n'(t) dt$
La suite de fonctions $(f_n')$ converge uniformément vers la fonction limite $g$. Par le théorème de continuité des limites uniformes de suites de fonctions continues (les dérivées $f_n'$ sont continues car $f_n$ est $C^1$), la limite $g$ est elle-même une fonction continue sur $[a,b]$.
Par conséquent, $g$ est Riemann-intégrable sur $[a,b]$.
La convergence uniforme sur $[a,b]$ de $(f_n')$ vers $g$ nous autorise formellement à appliquer le théorème d'interversion de la limite et de l'intégrale sur le segment de bornes $x_0$ et $x$ (qui est inclus dans $[a,b]$).
C'est-à-dire que pour tout $x \in [a,b]$ :
$\lim_{n \to \infty} \int_{x_0}^x f_n'(t) dt = \int_{x_0}^x \lim_{n \to \infty} f_n'(t) dt = \int_{x_0}^x g(t) dt$
De plus, par hypothèse de l'énoncé, la suite de nombres réels $(f_n(x_0))$ converge. Appelons $C$ sa limite : $\lim_{n \to \infty} f_n(x_0) = C$.
En passant à la limite (simple, pour chaque $x$ fixé) dans l'équation intégrale initiale, nous obtenons :
$\lim_{n \to \infty} f_n(x) = \lim_{n \to \infty} f_n(x_0) + \lim_{n \to \infty} \int_{x_0}^x f_n'(t) dt = C + \int_{x_0}^x g(t) dt$
Posons alors la fonction limite $f$ définie sur $[a,b]$ par :
$f(x) = C + \int_{x_0}^x g(t) dt$
La fonction $f$ est ainsi la limite simple de la suite $(f_n)$.
La fonction $g$ étant continue, son intégrale de Riemann $\int_{x_0}^x g(t) dt$ est, par le Théorème Fondamental de l'Analyse, l'unique primitive de $g$ qui s'annule en $x_0$.
Cela implique immédiatement que $f$ est une fonction différentiable sur $[a,b]$ et que sa dérivée est $f'(x) = g(x)$.

Il reste à démontrer que la convergence de $(f_n)$ vers $f$ est non seulement simple, mais bien uniforme.
Calculons l'écart entre $f_n(x)$ et $f(x)$ en utilisant leurs formes intégrales respectives :
$f_n(x) - f(x) = (f_n(x_0) - C) + \int_{x_0}^x (f_n'(t) - g(t)) dt$
En appliquant l'inégalité triangulaire (la valeur absolue d'une somme est inférieure ou égale à la somme des valeurs absolues, et la valeur absolue d'une intégrale est inférieure ou égale à l'intégrale de la valeur absolue) :
$|f_n(x) - f(x)| \leq |f_n(x_0) - C| + \left| \int_{x_0}^x |f_n'(t) - g(t)| dt \right|$
Soit $\epsilon > 0$.
Par la convergence uniforme de $(f_n')$ vers $g$ sur $[a,b]$, il existe un rang $N_1$ tel que pour tout $n \geq N_1$ et tout $t \in [a,b]$, $|f_n'(t) - g(t)| \leq \frac{\epsilon}{2(b-a)}$.
Par la convergence numérique de $(f_n(x_0))$ vers $C$, il existe un rang $N_2$ tel que pour tout $n \geq N_2$, $|f_n(x_0) - C| \leq \frac{\epsilon}{2}$.
Posons $N = \max(N_1, N_2)$. Pour tout $n \geq N$ et pour tout $x \in [a,b]$, majorons l'intégrale :
$\left| \int_{x_0}^x |f_n'(t) - g(t)| dt \right| \leq \left| \int_{x_0}^x \frac{\epsilon}{2(b-a)} dt \right| = \frac{\epsilon}{2(b-a)} |x - x_0|$
Puisque $x$ et $x_0$ appartiennent au segment $[a,b]$, la distance maximale entre eux est $|x - x_0| \leq b - a$.
Donc, l'intégrale est majorée par $\frac{\epsilon}{2(b-a)} (b-a) = \frac{\epsilon}{2}$.
En combinant les deux majorations pour l'inégalité globale :
$|f_n(x) - f(x)| \leq \frac{\epsilon}{2} + \frac{\epsilon}{2} = \epsilon$.
Cette inégalité finale ne dépend d'aucun $x$ particulier. En passant au supremum sur le compact $[a,b]$ :
$\sup_{x \in [a,b]} |f_n(x) - f(x)| \leq \epsilon$.
Ceci étant vrai pour tout $\epsilon > 0$ et pour un $n$ suffisamment grand, cela constitue la définition exacte de la limite : $\lim_{n \to \infty} \|f_n - f\|_{\infty} = 0$. La convergence de $(f_n)$ vers $f$ est donc rigoureusement uniforme sur $[a,b]$.

\newpage
\chapter*{Travaux Pratiques}
\addcontentsline{toc}{chapter}{Travaux Pratiques}


\chapter*{TP 1 : Visualisation de la convergence simple}

\textbf{Objectif :} Implémenter une fonction Python pour évaluer $f_n(x) = x^n$ sur $[0,1]$ et vérifier la convergence simple point par point via des assertions pour différentes valeurs de $x$.

\section*{Implémentation Pure Python (Validation Mathématique)}

\begin{lstlisting}[language=Python]
def f_n(n, x):
    return x ** n

def limite_f(x):
    return 0.0 if x < 1.0 else 1.0

# Validation par assertions
assert abs(f_n(100, 0.5) - limite_f(0.5)) < 1e-5
assert abs(f_n(10, 1.0) - limite_f(1.0)) < 1e-5
print("Convergence simple vérifiée aux points tests.")
\end{lstlisting}

\newpage


\chapter*{TP 2 : Calcul de la norme infinie}

\textbf{Objectif :} Créer un algorithme from scratch pour estimer $\sup_{x \in [0,1]} |f_n(x) - f(x)|$ en échantillonnant finement l'intervalle, afin d'observer la non-convergence uniforme de $x^n$.

\section*{Implémentation Pure Python (Validation Mathématique)}

\begin{lstlisting}[language=Python]
def f_n(n, x):
    return x ** n

def limite_f(x):
    return 0.0 if x < 1.0 else 1.0

def norme_infinie(n, points):
    erreurs = [abs(f_n(n, x) - limite_f(x)) for x in points]
    max_err = erreurs[0]
    for e in erreurs:
        if e > max_err: max_err = e
    return max_err

points = [i/1000.0 for i in range(1001)]
n_test = 100
err_max = norme_infinie(n_test, points)

# L'erreur reste proche de 1 car pour x proche de 1, x^n approche 1 tandis que f(x)=0
assert err_max > 0.9
print(f"L'erreur maximale pour n={n_test} est {err_max}, démontrant la NON-convergence uniforme.")
\end{lstlisting}

\newpage


\chapter*{TP 3 : Convergence uniforme de $x/n$}

\textbf{Objectif :} Vérifier empiriquement que $f_n(x) = x/n$ converge uniformément sur $[-10, 10]$ en calculant le max de l'erreur, et vérifier par assertions que cette erreur maximale tend vers $0$.

\section*{Implémentation Pure Python (Validation Mathématique)}

\begin{lstlisting}[language=Python]
def f_n(n, x):
    return x / n

def norme_infinie_segment(n, a, b, num_points=1000):
    pas = (b - a) / num_points
    points = [a + i*pas for i in range(num_points + 1)]
    erreurs = [abs(f_n(n, x)) for x in points] # f(x) = 0
    max_err = erreurs[0]
    for e in erreurs:
        if e > max_err: max_err = e
    return max_err

err_10 = norme_infinie_segment(10, -10, 10)
err_100 = norme_infinie_segment(100, -10, 10)
err_1000 = norme_infinie_segment(1000, -10, 10)

assert err_10 > err_100 > err_1000
assert err_1000 <= 10 / 1000
print("La décroissance de la norme infinie vers 0 est vérifiée.")
\end{lstlisting}

\newpage


\chapter*{TP 4 : Le phénomène de Gibbs discret}

\textbf{Objectif :} Modéliser une série de Fourier tronquée d'un signal créneau en Python pur et observer l'évolution de la norme infinie de l'erreur, illustrant l'absence de convergence uniforme près des discontinuités.

\section*{Implémentation Pure Python (Validation Mathématique)}

\begin{lstlisting}[language=Python]
import math

def fourier_partiel(n, x):
    somme = 0
    for k in range(1, n+1, 2):
        somme += (4 / (math.pi * k)) * math.sin(k * x)
    return somme

# f(x) = 1 sur ]0, pi[, -1 sur ]-pi, 0[
def norme_inf_erreur(n, num_points=500):
    pas = math.pi / num_points
    # On évalue sur ]0, pi[ pour éviter les bornes
    points = [pas + i*pas for i in range(num_points - 1)]
    erreurs = [abs(fourier_partiel(n, x) - 1.0) for x in points]
    max_err = erreurs[0]
    for e in erreurs:
        if e > max_err: max_err = e
    return max_err

err_10 = norme_inf_erreur(10)
err_50 = norme_inf_erreur(50)

# Le pic d'erreur reste autour de 18% (phénomène de Gibbs)
assert err_10 > 0.15 and err_50 > 0.15
print("Le phénomène de Gibbs empêche la convergence uniforme.")
\end{lstlisting}

\newpage


\chapter*{TP 5 : Validation du théorème de Dini}

\textbf{Objectif :} Implémenter une suite de fonctions continues décroissant vers une limite continue sur un maillage discret. Écrire un script pour trouver le $N$ tel que pour tout $n \geq N$, l'erreur maximale soit inférieure à $\epsilon$, validant la convergence uniforme.

\section*{Implémentation Pure Python (Validation Mathématique)}

\begin{lstlisting}[language=Python]
def f_n(n, x):
    return 1 / (1 + n * x)

# Sur [a, 1] avec a > 0, la suite converge uniformément vers 0.
# Vérifions que pour eps > 0 donné, on peut trouver un N global.
def trouver_N_uniforme(eps, a, b, num_points=1000):
    pas = (b - a) / num_points
    points = [a + i*pas for i in range(num_points + 1)]

    n = 1
    while True:
        max_err = 0
        for x in points:
            err = abs(f_n(n, x))
            if err > max_err: max_err = err

        if max_err < eps:
            return n
        n += 1

a = 0.1
b = 1.0
eps = 0.05
N_requis = trouver_N_uniforme(eps, a, b)

# Validation : Pour n >= N_requis, max_err doit être < eps
max_err_validation = 0
pas = (b - a) / 100
for i in range(101):
    x = a + i*pas
    err = abs(f_n(N_requis, x))
    if err > max_err_validation: max_err_validation = err
assert max_err_validation < eps
print(f"Théorème de Dini validé : N trouvé = {N_requis}")
\end{lstlisting}

\newpage
\end{document}